%
% Copyright 2016 Erint Labs OÜ
%
% Licensed under the EUPL, Version 1.1 or – as soon they
% will be approved by the European Commission - subsequent
% versions of the EUPL (the "Licence");
% You may not use this work except in compliance with the
% Licence.
% You may obtain a copy of the Licence at:
%
% https://joinup.ec.europa.eu/software/page/eupl
%
% Unless required by applicable law or agreed to in
% writing, software distributed under the Licence is
% distributed on an "AS IS" basis,
% WITHOUT WARRANTIES OR CONDITIONS OF ANY KIND, either
% express or implied.
% See the Licence for the specific language governing
% permissions and limitations under the Licence.
%
%
% Copyright 2016 Erint Labs OÜ
%
% Licensed under the EUPL, Version 1.1 or – as soon they
% will be approved by the European Commission - subsequent
% versions of the EUPL (the "Licence");
% You may not use this work except in compliance with the
% Licence.
% You may obtain a copy of the Licence at:
%
% https://joinup.ec.europa.eu/software/page/eupl
%
% Unless required by applicable law or agreed to in
% writing, software distributed under the Licence is
% distributed on an "AS IS" basis,
% WITHOUT WARRANTIES OR CONDITIONS OF ANY KIND, either
% express or implied.
% See the Licence for the specific language governing
% permissions and limitations under the Licence.
%
\documentclass[a4paper]{article}
\usepackage[landscape,top=1cm, bottom=1.5cm, left=1cm, right=1cm]{geometry}
% TODO: non XeLaTeX compliation support
\usepackage[colorlinks]{hyperref}
\usepackage{paracol}
\usepackage{fancyref}
\usepackage{fontspec}
\usepackage{xltxtra}
\usepackage{ifthen}
\usepackage{multirow}
\usepackage{xunicode}
\usepackage{indentfirst}
\usepackage{fancyhdr}
\XeTeXinputencoding utf8
\usepackage{polyglossia}
\setdefaultlanguage{ukrainian}
\setotherlanguages{english,estonian}
\setromanfont[Mapping=tex-text]{Libertinus Serif}

% TODO: remove color borders maybe?
% TODO: better document class?
% TODO: make hyperref only use English names for TOC entries
% TODO: prevent letters in numeric fields
% TODO: stylish small capitals
% WARNING: current Estonian version may contain machine translation.
% DO NOT use it right now if you don't know what you're doing

% Footer to contain the template
\pagestyle{fancy}
\fancyhead{}
\renewcommand{\headrulewidth}{0pt}
\lfoot{\large\thepage}
\cfoot{}
\newcommand\docurl{https://github.com/Ecognize/contracts/blob/master/\jobname.tex}
\rfoot{\scriptsize Шаблон доступний за адресою / Template available at / Šabloon on kättesaadav aadressil : \href{\docurl}{\docurl} }

% Entry boxes with \raisebox
\newcommand\field[3]{\raisebox{-#2pt}{\TextField[borderwidth=0,#3,name=#1]{}}}
\newcommand\fieldrws[4]{\field{#1}{#2}{width=#3pt,charsize=#4pt}}
\newcommand\fieldhs[3]{\field{#1}{2}{height=#2pt,charsize=#3pt}}
\newcommand\fieldwhs[4]{\field{#1}{2}{width=#2pt,height=#3pt,charsize=#4pt}}
\newcommand\fieldw[2]{\field{#1}{2}{width=#2pt}}
\newcommand\fieldtw[2]{\field{#1}{2.2}{height=10pt,charsize=8pt,width=#2pt}}

% Most typical variations
\newcommand\fieldline[1]{\fieldw{#1}{250}}
\newcommand\fieldterm[1]{\fieldhs{#1}{12}{10}}
\newcommand\fieldt[1]{\fieldtw{#1}{40}}
\newcommand\fieldb[1]{\fieldwhs{#1}{150}{14}{10}}
\newcommand\fieldbs[1]{\fieldwhs{#1}{150}{12}{8}}

% Three column text in UK, EN, and ET
\newcommand\freetextcommon[3]{\begin{ukrainian}#1\end{ukrainian}\switchcolumn\begin{english}#2\end{english}\switchcolumn\begin{estonian}#3\end{estonian}}
\newcommand\freetextnoalign[3]{\freetextcommon{#1}{#2}{#3}\swtichcolumn}
\newcommand\freetext[3]{\freetextcommon{#1}{#2}{#3}\switchcolumn*}
\newcommand\freetextnoindent[3]{\freetext{\noindent#1}{\noindent#2}{\noindent#3}}
\newcommand\multibreak{\freetext{\pagebreak}{\pagebreak}{\pagebreak}}
% Same with section header
\newcommand\clause[6]{\freetext{\section{#1}#4}{\section{#2}#5}{\section{#3}#6}}
\newcommand\clausenewpage[6]{\freetext{\pagebreak\section{#1}#4}{\pagebreak\section{#2}#5}{\pagebreak\section{#3}#6}}
% Cute symbols
\renewcommand\thesection{§\arabic{section}}
\renewcommand\thesubsection{§\arabic{section}.\arabic{subsection}}
% Hyperlinks to Appendixes and sections
\newcommand\smartref[2]{\hyperref[#1]{#2\ref{#1}}}
% This is lame, but I have no will to make it better
\newcommand{\myloop}[3]{\stepcounter{#1}#3\ifthenelse{\value{#1} < \value{#2}}{\\\myloop{#1}{#2}{#3}}{\\}}


\begin{document}
  % We are stylish
  % TODO: generalise
  \backgroundsetup{contents={\includegraphics[width=0.3cm]{ecognize.pdf}},angle=0,position={0,0}}


  \begin{Form}
    \title{Угода про надання послуг / Service provision agreement / Töövõtuleping\\№ \fieldrws{contractno}{6}{250}{14}}
    % TODO: using author field for place is ugly, let's use another docuement class in future
    \author{\fieldterm{placeUK} / \fieldterm{placeEN} / \fieldterm{placeET}}
    \date{\today / \textenglish{\today} / \textestonian{\today}}
    \maketitle
    \thispagestyle{fancy}

    \begin{paracol}{3}
      \freetext % TODO do we have to support people who refused an ID number btw? Add this if necessary
        {Нинішьною \textbf{\fauxsc{угодою}}\\
          \fieldline{customernameUK},\\
          якого(-у) представляє (на підставі)\\
          \fieldline{customerpersonUK}\\
          (надалі \textbf{\fauxsc{замовник}}) з однієї сторони, а з іншої\\
          \fieldline{providernameUK},\\
          якого(-у) представляє (на підставі)\\
          \fieldline{providerpersonUK}\\
          (надалі \textbf{\fauxsc{виконавець}}, а разом із \fauxsc{замовником} надалі \textbf{\fauxsc{сторони}}) домовились про наступне:
        }
        { With this \textbf{\textsc{agreement}}\\
          \fieldline{customernameEN}\\
          represented by (on the basis of)\\
          \fieldline{customerpersonEN}\\
          (referred to as \textbf{\textsc{customer}}) from one side and\\
          \fieldline{providernameEN}\\
          represented by (on the basis of)\\
          \fieldline{providerpersonENET}\\
          (referred to as \textbf{\textsc{provider}}, referred to collectively with the \textsc{customer} as the \textbf{\textsc{parties}}) have agreed on the following:
          }
        {Lorem ipsum dolor sit amet, consectetur adipiscing elit. Pellentesque arcu leo, malesuada vitae eleifend vitae, ultrices quis nulla. Curabitur eget pharetra tortor. Phasellus nisi massa, facilisis sed dui vestibulum, maximus malesuada mauris. Duis dapibus consectetur tellus, a ornare lacus porta a. Nulla dignissim velit in lectus tincidunt, nec vulputate eros malesuada. Aliquam aliquam luctus maximus. Proin facilisis blandit turpis, eget sodales purus efficitur ut. Morbi vel faucibus dolor. In vel dapibus erat, in mollis nulla. Nullam id porttitor eros. Nullam vulputate dui in faucibus aliquet. Ut eu cursus orci. Nunc sed finibus lorem.}
      \clause
        {Мова та терміни}
        {Language and terms}
        {Eesti}
        {Нинішню \fauxsc{угоду} укладено у трьох версіях: українською, англійською та естонською мовами. Офіційні назви сторін \fauxsc{угоди} наведені у версії, що відповідає державній мові країни, в якій відповідна сторона є резидентом, а також неофційно адаптовані у інших версіях. У разі сумнівів щодо відповідності версій, українська вважатиметься оригіналом, окрім офіційних назв. Виключно з метою перекладу відповідність ор\-га\-ні\-за\-цій\-но-правових форм господарювання між мовами \fauxsc{угоду} наведено у \smartref{app:correspondence}{додатку }. \textbf{\fauxsc{третьою особою}} позначено будь-яку іншу сторону, що має договірні стосунки із \fauxsc{замовником}, але є відмінною від \fauxsc{виконавця}. В межах коммунікації щодо виконання \fauxsc{угоди} між сторонами можуть виникати усні та письмові \textbf{\fauxsc{домовленності}}, що не є текстовою частиною угоди і не потребують підпису \fauxsc{сторін}.
        }
        {This \textsc{agreement} is concluded in three langauge versions, na\-me\-ly: Ukrainian, English and Estonian. The official \textsc{agreement} party names are provided as they are in the version corresponding to the state language of the country where the respective party is resident in and are unofficially adapted in other translations. In case of doubts concerning the conformity between the versions, the Ukrainian one shall be considered as the original, except for the official names. For the purpose of translation only the correspondence between the types of business entities in \textsc{agreement}'s languages is provided in \smartref{app:correspondence}{appendix }. Any other party distinct from the \textsc{provider} that is contracted by the \textsc{customer} is referred to as the \textbf{\textsc{third party}}. During the communucation concerning the execution of the \textsc{agreement}, the \textsc{parties} may conclude separate oral or written \textbf{\textsc{arrangements}} which are separate from the text of the \textsc{agreement} and do not require being signed.
        }
        {Lorem ipsum dolor sit amet, consectetur adipiscing elit. Pellentesque arcu leo, malesuada vitae eleifend vitae, ultrices quis nulla. Curabitur eget pharetra tortor. Phasellus nisi massa, facilisis sed dui vestibulum, maximus malesuada mauris. Duis dapibus consectetur tellus, a ornare lacus porta a. Nulla dignissim velit in lectus tincidunt, nec vulputate eros malesuada. Aliquam aliquam luctus maximus. Proin facilisis blandit turpis, eget sodales purus efficitur ut. Morbi vel faucibus dolor. In vel dapibus erat, in mollis nulla. Nullam id porttitor eros. Nullam vulputate dui in faucibus aliquet. Ut eu cursus orci. Nunc sed finibus lorem.}
      \clause
        {Модальні дієслова}
        {Modal verb use}
        {Eesti}
        {Модальні дієслова української мови в межах \fauxsc{угоди} вживаються в наступному сенсі: \textbf{має} — обов’язок \fauxsc{сторони}, невиконання якого є порушенням \fauxsc{угоди}; \textbf{може} — право \fauxsc{сторони}, виконання або невиконання якого не є порушенням \fauxsc{угоди}; \textbf{не має} — заборона \fauxsc{стороні}, порушення якої є порушенням \fauxsc{угоди}.}
        {For the scope of the \textsc{agreement} the following interpretation of the English modal verbs is assumed: \textbf{shall} — a duty of the \textsc{party} failure to provide which constitutes a breach of the \textsc{agreement}; \textbf{may} — a right of the \textsc{party} the provision or the failure to provide which does not constitute a breach of the \textsc{agreement}; \textbf{shall not} — a prohibition for the \textsc{party}, the breach of which is the breach of the \textsc{agreement}.}
        {Lorem ipsum dolor sit amet, consectetur adipiscing elit. Pellentesque arcu leo, malesuada vitae eleifend vitae, ultrices quis nulla. Curabitur eget pharetra tortor. Phasellus nisi massa, facilisis sed dui vestibulum, maximus malesuada mauris. Duis dapibus consectetur tellus, a ornare lacus porta a.}
      \clause
        {Предмет}
        {Object}
        {Eesti}
        {С\fauxsc{торони} домовляються про те, що \fauxsc{виконавець} \textit{має} виконувати, а \fauxsc{замовник} \textit{має} приймати та у разі \emph{належного виконання} оплачуватиме послуги, зазначені у \smartref{app:services}{додатку }. Послуги \textit{можуть} виконуватимуться на користь \fauxsc{замовника} безпосередьно, або на користь \fauxsc{третьої сторони}, із якою \fauxsc{замовник} має відповідну \fauxsc{угоду} або \fauxsc{домовленність}.

        \subsection{Належне виконання}
        Під \emph{належним виконанням} \fauxsc{сторони} розуміють виконання послуги відповідно до технічного опису завдання, яке \fauxsc{замовник} передає \fauxsc{виконавцеві} у усній, матеріальній або електронній формі, що було здійснено не пізніше визначеного крайнього терміну. Якщо такий термін не визначений прямо, \fauxsc{сторони} \emph{мають} вважати один календарний місяць з часу отрмання технічного опису завдання крайнім терміном за замовченням.
        }
        {The \textsc{parties} agree that the \textsc{provider} \textit{shall} provide while the \textsc{customer} \textit{shall} accept and pay in the case of the the diligent provision of the services indicated in \smartref{app:services}{appendix }. The services \textit{may} be provided in favour of the \textsc{customer} directly or to a \textsc{third party} with which the \textsc{customer} has a respective \textsc{agreement} or \textsc{arrangement}.

        \subsection{Diligent priovision}
        The \emph{diligent provision} is understood by the \textsc{parties} as the execution of service according to the technical task description that the \textsc{customer} supplies to the \textsc{provider} in oral, material or electronic form that was complete before the deadline specified. In case the term is not explicitly defined, the \textsc{parties} \emph{shall} consider one calendar month from the date when the technical description was recevied to be a default deadline.
        }
        {Lorem ipsum dolor sit amet, consectetur adipiscing elit. Pellentesque arcu leo, malesuada vitae eleifend vitae, ultrices quis nulla. Curabitur eget pharetra tortor. Phasellus nisi massa, facilisis sed dui vestibulum, maximus malesuada mauris. Duis dapibus consectetur tellus, a ornare lacus porta a.}
      \clase
        {Ціна}
        {Price}
        {Eesti}
        {Одиницею виміру наданих послуг є година часу, витраченого на їх надання. Ціна однієї години зазначена у \smartref{app:services}{додатку }. Відсутність ціни на послугу, або її нульове значення означає згоду на безоплатне виконання відповідної послуги \fauxsc{виконавцем}.}
        {An hour of time spent for the service provision is to be considered the unit of services provided. \smartref{app:services}{Appendix } determines the hourly price. The lack of the price for the service or its zero value means that \textsc{provider} agrees to supply the respective service free of charge.}
        {}
      \clause
        {Кількість}
        {Quantity}
        {Eesti}
        {Кількість послуг в годинах визначається \fauxsc{виконавцем} відповідно до складності та конкретики завдань. Технічний опис завдань \emph{може} містити максимальну кількість годин, які \fauxsc{замовник} \emph{має} оплатити, з чим \fauxsc{виконавець} погоджується.}
        {The hourly quantity of services provided is determined by the \textsc{provider} according to the complexity and the pecularities of the tasks. A task technical description \emph{may} contain an upper limit of hours that the \textsc{customer} \emph{shall} pay to which the \textsc{provider} agrees.}
        {Lorem ipsum dolor sit amet, consectetur adipiscing elit. Pellentesque arcu leo, malesuada vitae eleifend vitae, ultrices quis nulla. Curabitur eget pharetra tortor. Phasellus nisi massa, facilisis sed dui vestibulum, maximus malesuada mauris. Duis dapibus consectetur tellus, a ornare lacus porta a.}
      \clause
        {Вартість}
        {Price}
        {Eesti}
        {Вартість послуг, здійснених у межах \fauxsc{угоди} \textit{не має} перевищувати \fieldw{fullprice}{75} €. У випадку потреби здійснення послуг на суму, більше цього обмеження, \fauxsc{сторони} \textit{мають} укласти нову \fauxsc{угоду}.}
        {The price of the services provided within the \textsc{agreement}'s scope \textit{shall not} exceed \fieldw{fullprice}{75} €. Should the need of the service provision above this sum arise, the \textsc{parties} \textit{shall} conclude a new \textsc{agreement}.}
        {Lorem ipsum dolor sit amet, consectetur adipiscing elit. Pellentesque arcu leo, malesuada vitae eleifend vitae, ultrices quis nulla. Curabitur eget pharetra tortor.}
      \clause % TODO check EU directive
        {Платежі}
        {Payments}
        {Eesti}
        {\label{sec:payment}Послуги \textit{мають} виконуватись за умов повної післясплати. \fauxsc{виконавець} \textit{має} вести облік часу виконання \textit{має} складати рахунки-фактури на цій підставі. Сплата \fauxsc{замовником} рахунку фактури \textit{має} визнаватись фактом прийняття виконання \fauxsc{замовником} послуг \fauxsc{виконавця} та відсутності претензій до їх якості. З\fauxsc{амовник} оплачує послуги шляхом міжнародного банківського переказу на валютний рахунок \fauxsc{виконавця}.

        \subsection{Формат рахунку-фактури}
        \begin{enumerate}
          \item Рахунок-фактура \textit{має} виставлятись у електронній формі та містити електронно-цифровий підпис \fauxsc{виконавця}.
          \item Рахунок-фактура \textit{має} відповідати вимогам \href{http://bank.gov.ua/doccatalog/document?id=19208488}{повідомлення НБУ №22-01012/46746 від 7 липня 2015 р.}.
          \item Рахунок-фактура \textit{може} мати формат сумісний із вимогами \href{http://eur-lex.europa.eu/legal-content/EN/TXT/?uri=CELEX:32014L0055}{директиви ЄУ 2014/55/EU від 13 жовтня 2014 р.}
        \end{enumerate}
        }
        {The services \textit{shall} be provided on the condition of complete post-payment. The \textsc{provider} \textit{shall} keep track of the time spent for the service provision and \textit{shall} issue invoices based on it. The payment for the invoice by the \textsc{customer} \textit{shall} be considered as the acceptance of the services executed by the \textsc{provider} and the lack of claims concerning their quality. The \textsc{customer} pays the services via an international bank transfer on the \textsc{provider}'s account.

        \subsection{Invoice format}
        \begin{enumerate}
          \item The invoice \textit{shall} be issued in an electronic form and digitally signed by the \textsc{provider}.
          \item The content of the invoice \textit{shall} correspond to the requirements of \href{http://bank.gov.ua/doccatalog/document?id=19208488}{the NBoU notice №22-01012/46746 dated July 7, 2015}.
          \item The invoice \textit{may} have a format compatible with \href{http://eur-lex.europa.eu/legal-content/EN/TXT/?uri=CELEX:32014L0055}{the EU directive 2014/55/EU dated October 13, 2014.}
        \end{enumerate}
        }
        {Lorem ipsum dolor sit amet, consectetur adipiscing elit. Pellentesque arcu leo, malesuada vitae eleifend vitae, ultrices quis nulla. Curabitur eget pharetra tortor. Phasellus nisi massa, facilisis sed dui vestibulum, maximus malesuada mauris. Duis dapibus consectetur tellus, a ornare lacus porta a. Nulla dignissim velit in lectus tincidunt, nec vulputate eros malesuada. Aliquam aliquam luctus maximus. Proin facilisis blandit turpis, eget sodales purus efficitur ut. Morbi vel faucibus dolor. In vel dapibus erat, in mollis nulla. Nullam id porttitor eros. Nullam vulputate dui in faucibus aliquet. Ut eu cursus orci. Nunc sed finibus lorem.}
      \clause
        {Місцевість}
        {Locality}
        {Eesti}
        {С\fauxsc{торони} не мають обмежень щодо місцевості виконання послуг. В\fauxsc{иконавець} \textit{може} виконувати послуги віддаленно, комунікувати із \fauxsc{замовником} та передавати результати виконаних послуг завдяки мережі Інтернет або іншим прийнятним для \fauxsc{сторін} способом. За \fauxsc{домовленністю} \fauxsc{замовник} \textit{може} тимчасово запрошувати \fauxsc{виконавця} до виконання послуг у визначеній \fauxsc{замовником} місцевості.

        \subsection{Витрати на подорож}
        З\fauxsc{амовник} \textit{може} оплачувати витрати \fauxsc{виконавця} на подорожі, пов'язані із виконанням \fauxsc{угоди} як власні витрати. % TODO: ask Estonian tax office about that
        }
        {The \textsc{parties} don't put forth any limitations on the locality of the service provision. The \textsc{provider} \textit{may} provide the services remotely, communicating and delivering the results of the service provision to the \textsc{customer} via the Internet or via another mutually acceptable means. Given an \textsc{arrangement}, the \textsc{customer} \textit{may} temporarily invite the \textsc{provider} to supply services in the locality specified by the \textsc{customer}.

        \subsection{Travel Expenses}
        The \textsc{customer} \textit{may} pay for the travelling expenses suffered by the \fauxsc{provider} related to the execution of the \textsc{agreement} as their own expense.
        }
        {Lorem ipsum dolor sit amet, consectetur adipiscing elit. Pellentesque arcu leo, malesuada vitae eleifend vitae, ultrices quis nulla. Curabitur eget pharetra tortor. Phasellus nisi massa, facilisis sed dui vestibulum, maximus malesuada mauris. Duis dapibus consectetur tellus, a ornare lacus porta a. Nulla dignissim velit in lectus tincidunt, nec vulputate eros malesuada. }
      \clause
        {Обладнання}
        {Equipment}
        {Eesti}
        {В\fauxsc{иконавець} \textit{має} забезпечити собі необхідне для виконання послуг обладнання та інфраструктуру власним коштом. Витрати, пов’язані із підтриманням робочого середовища \textit{мають} бути включені \fauxsc{виконавцем} у вартість послуг. 

        \subsection{Незамінне обладнання}
        З\fauxsc{амовник} \textit{може} надавати \fauxsc{виконавцю} у тичасове користування свою власність, таку як прототипи, продукти, унікальне обладнання \textit{тощо} якщо послуги виконуються безпосередньо із нею, в разі чого \fauxsc{замовник} \textit{має} оплачувати транспортні витрати.}
        {The \textsc{provider} \textit{shall} provide themself with the equipment and infrastructure needed for service provision at their own expense. The expenses arsising from the maintenance of the work environment \textit{shall} be included into the service price by the \textsc{provider}. 

        \subsection{Non-substitutable equipment}
        Whereas the services are to be peformed directly on the \textsc{customer}'s property such as unique equipment, products, prototypes, \textit{etc} the \textsc{customer} \textit{may} temporarily provide it to the \textsc{provider}. In such case the \textsc{customer} \textit{shall} pay the transportation expenses.}
        {Lorem ipsum dolor sit amet, consectetur adipiscing elit. Pellentesque arcu leo, malesuada vitae eleifend vitae, ultrices quis nulla. Curabitur eget pharetra tortor. Phasellus nisi massa, facilisis sed dui vestibulum, maximus malesuada mauris. Duis dapibus consectetur tellus, a ornare lacus porta a. Nulla dignissim velit in lectus tincidunt, nec vulputate eros malesuada. }
      \clause
        {Стосунки сторін}
        {Parties's relations}
        {Eesti}
        {С\fauxsc{торони} зазначають, що \fauxsc{виконавець}:
        \begin{enumerate}
          \item Надає послуги за завданням самостійно, під власну відповідальність та ризик.
          \item Не підпорядковується штатному розпорядку \fauxsc{замовника}.
          \item Не користується звичайним обладнанням, приміщеннями та транспортом \fauxsc{замовника}.
          \item Переважно обирає час там місце виконання послуг на власний розсуд.
          \item Не є резидентом Республіки Естонія.
          \item Не користується оплачуваною відпусткою та іншими пільгами.
        \end{enumerate}

        Відповідно до наведеного, \fauxsc{сторони} вважають, що \fauxsc{виконавець} не є найманим працівником \fauxsc{замовника} у розумінні \href{http://zakon2.rada.gov.ua/laws/show/322-08}{Кодексу законів про працю України}, а також \href{http://zakon24.ee/zakon-o-trudovom-dogovore/}{Закону про трудову угоду Республіки Естонія}.

        \subsection{Підприємницьтво}
        Відповідно до \href{http://zakon3.rada.gov.ua/laws/show/959-12/parao138\#o138}{статті 5 Закону України «Про зовнішньоекономічну діяльність»}, \fauxsc{виконавець}, який не є юридичною особою, \textit{має} бути зареєстрований як підприємець.

        \subsection{Право представлення}
        Під час дії \fauxsc{угоди} \fauxsc{виконавець} \textit{може} асоціювати себе із брендом та ім'ям \fauxsc{замовника}, користуватись його атрибутикою під час участі у урочистих заходах, конференціях тощо, за \fauxsc{домовленністю}.

        \subsection{Пов'язані особи}
        В\fauxsc{иконавець} або його представник \textit{може} володіти долею у підприємстві \fauxsc{замовника} достатньою за розміром для його визначення як пов’язаної особи. З огляду на це \fauxsc{сторони} домовляються що вони \textit{мають} за потреби дотримуватись вимог щодо документації транзакцій трансфертного ціноутворення країн де вони є резидентами. Положення \fauxsc{угоди} \textit{мають} бути виконані за принципом витягнутої руки та відсутності конфлікту інтересів. У разі виявлення невідповідності \textit{угоди} наведеним вимогам, \fauxsc{сторони} \textit{мають} невідкладно розірвати цю \fauxsc{угоду} та укласти нову із відпвідними виправленнями.}
        {The \textsc{parties} note that the \textsc{provider}:
        \begin{enumerate}
          \item Provides the services independently at his own risk and responsibility.
          \item Is not subject to the working schedule of the \textsc{customer}.
          \item Does not make use of the common equipment, premises or transport of the \textsc{custonmer}.
          \item Largely uses their own reasoning deciding the place and the time of the service provision.
          \item Does not reside in Republic of Estonia.
          \item Does not make use of paid vacation and other bonuses.
        \end{enumerate}

        Considering the aforementioned, the \textsc{parties} assume that the \textsc{provider} is not deemed to be an employee of the \textsc{customer} in the definitions of both \href{http://zakon2.rada.gov.ua/laws/show/322-08}{Labour Code of Ukraine} and \href{https://www.riigiteataja.ee/en/eli/530102013061/consolide}{Employment Contracts Act of Estonia}.

        \subsection{Entrepreneurship}
        In accordance to \href{http://zakon3.rada.gov.ua/laws/show/959-12/parao138\#o138}{the article 5 of Ukrainian Law on External Economic Activity} the \textsc{provider} which is not a legal person \textit{shall} be registered as an entrepreneur.

        \subsection{Right to represent}
        The \textsc{provder}, \textit{may} associate themselves with the \textsc{customer}'s name and brand, make use of \textsc{customer}'s symbolics for the purpose of the ceremonial events, conferences and so on, based on an \textsc{arrangement}.

        \subsection{Associated entities}
        The \textsc{provider} or their representative \textit{may} own a share of the \textsc{customer}'s company big enough for the \textsc{provider} and the \textsc{customer} to be considered as the associated entities. In such case the \textsc{parties} agree that they \textit{shall} document the transactions in accordance to the requirements for the transfer pricing set by the countries they are resident in as needed. The clauses of the \textsc{agreement} \textit{shall} be drafted under the principle of arm's length and the lack of conflict of interest. Should the \textsc{agreement} be found not to comply with the aforementioned conditions, the \textsc{parties} \textsc{shall} immediately dissolve this \textsc{agreement} and conclude another one with respective corrections.
        }
        {Lorem ipsum dolor sit amet, consectetur adipiscing elit. Pellentesque arcu leo, malesuada vitae eleifend vitae, ultrices quis nulla. Curabitur eget pharetra tortor. Phasellus nisi massa, facilisis sed dui vestibulum, maximus malesuada mauris. Duis dapibus consectetur tellus, a ornare lacus porta a. Nulla dignissim velit in lectus tincidunt, nec vulputate eros malesuada. Aliquam aliquam luctus maximus. Proin facilisis blandit turpis, eget sodales purus efficitur ut. Morbi vel faucibus dolor. In vel dapibus erat, in mollis nulla. Nullam id porttitor eros. Nullam vulputate dui in faucibus aliquet. Ut eu cursus orci. Nunc sed finibus lorem.}
      \clause
        {Податки та облік}
        {Taxation and accounting}
        {Eesti}
        {С\fauxsc{торони} \textit{мають} самостійно вести бухгалтерський облік та сплату податків незалежно одне від одного, відповідно до вимог країни, де вони є резидентами. В\fauxsc{иконавець} \textit{має} включати витрати, пов’язані із сплатою власних податків та зборів у ціну послуг.

        \subsection{Податок на додану вартість}
        Якщо імпорт послуг \fauxsc{виконавця} виникає зобов'язання щодо сплати податку на додану вартість, \fauxsc{замовник} \textit{має} вести відповідний облік та звітність відповідно до принципу зворотнього оподаткування.
        }
        {The \textsc{parties} \textit{shall} do their accounting and pay their taxes on their own independently from each other given the rules of the country they are resident in. The \textsc{provider} \textit{shall} include the expenses related to the payment of his taxes to the service price.

        \subsection{Value added tax}
        If the services imported from the \textsc{provider} cause value added tax to be charged, the \textsc{provider} \textit{shall} conduct respective accounting and reporting in accoradnce to the reverse charge principle.
        }
        {Lorem ipsum dolor sit amet, consectetur adipiscing elit. Pellentesque arcu leo, malesuada vitae eleifend vitae, ultrices quis nulla. Curabitur eget pharetra tortor. Phasellus nisi massa, facilisis sed dui vestibulum, maximus malesuada mauris. Duis dapibus consectetur tellus, a ornare lacus porta a. Nulla dignissim velit in lectus tincidunt, nec vulputate eros malesuada. }
      \clause
        {Майнові права}
        {Property rights}
        {}
        {
        Обладнання, яке передане \fauxsc{виконавцю} для виконання послуг залишається у власності \fauxsc{виконавця}. У разі створення матеріальних об'єктів у ході виконання послуг, вони \textit{мають} передаватись у власність \fauxsc{замовника} завдяки окремим \fauxsc{умовам} або \fauxsc{домовленностям}.

        \subsection{Інтелектуальна власність}
        В\fauxsc{иконавець} \textit{має} передавати \fauxsc{замовнику} майнові авторські права на результати виконання послуг згідно до \fauxsc{угоди}, такі як програмний код, креслення, звіти, моделі, графічні матеріали, дані досліджень, патенти, документація тощо, але не обмежуючись цим переліком. В\fauxsc{иконавець} втім матиме право вважатись автором цього матеріалу та використовувати його в цілях портфоліо, якщо це не суперечить додатковим \fauxsc{домовленностям} про нерозголошення. У випадку, якщо \fauxsc{виконавець} виконує послуги на користь \fauxsc{третьої особи} від імені \fauxsc{замовника}, домовленість \fauxsc{замовника} та \fauxsc{третьої особи} щодо порядку передачі авторських прав матиме приорітет над положеннями цієї угоди.}
        {
        The equipment transferred to the \textsc{provider} for the service provision stays in the \textsc{customer}'s property. Should the material objects be the results of the service provision, they \textit{shall} be transferred to the \textsc{customer}'s property using separate \textsc{agreements} or \textsc{arrangements}.

        \subsection{Intellectual property}
        The \textsc{provider} shall transfer the copyright over the results of the services provided in the scope of \textsc{agreement} including, but not limited to software source code, schematics, reports, models, graphical material, research data, patents, documentation \emph{et ce\-te\-ra} to the \textsc{customer}. The \textsc{provider} however shall be considered the author of the respective material and allowed to use it in the purposes of porfolio unless such act would contradcit additional non-disclosure \textsc{arrangements}. In cases when the \textsc{provider} provides the services to the \textsc{third party} on the \textsc{customer}'s behalf, the agreement on copyright transfer between the \textsc{third party} and the \textsc{customer} would take priority over the provisions of this agreement.}
        {}
      \clause
        {Форс-мажор}
        {Force majeure}
        {}
        {Невиконання \fauxsc{стороною} обов'язків згідно \fauxsc{угоди}, що було спричинено обставинами непереборної дії, такими як дії уряду, війни, стихійні лиха \emph{та інші}, не \emph{має} вважатись порушенням \fauxsc{угоди}.}
        {The failure of the \textsc{party} to comply with their obligations arising from the \textsc{agreement} that is caused by the force majeure such as actions of the governments, wars, natural disasters \emph{et cetera}, \emph{shall} not be considered to be a breach of the \textsc{agreement}.}
        {}
      \clause
        {Урегулювання спорів}
        {Dispute resolution}
        {}
        {У\fauxsc{года} укладена у відповідності до законодавства України та Естонської Республіки. С\faux{торони} визначають юрисдикцію суду повіту Гар'юмаа для вирішення конфліктних ситуацій.}
        {The \textsc{agreement} is concluded in accordance to the laws of Ukraine and Estonian Republic. The \textsc{parties} define the jurisdiction of Harju county court for the settlement of conflicts.}
        {}
      \clause
        {Набуття чинності}
        {Entry into force}
        {}
        {С\fauxsc{торони} підписують \fauxsc{угоду} та її похідні документи за допомогою електронно-цифрового підпису, засвідченого у відповідних країнах за умови, що такий підпис вважатиметься чинним країнами, в яких \fauxsc{сторони} є резидентами, а також у країні, призначеній для вирішення спорів між \fauxsc{сторонами} щодо \fauxsc{угоди} та її виконання.

        \subsection{Матеріальна копія}
        У разі технічної, або юридичної неможливості накласти на \fauxsc{угоду} електронно-цифровий підпис, \fauxsc{угода} виготовляється у паперовому форматі в двох примірниках та підписується на кожній сторінці ліворуч — \fauxsc{замовником}, а праворуч — \fauxsc{виконавцем}. У\fauxsc{года} набуває чинності із момента її підписання обома \fauxsc{сторонами}.}
        {The \textsc{parties} sign the \textsc{agrement} and its derivative documents by means of digital signature certified by respective countries on the condition that such a signature would be considered valid in the countries the \textsc{parties} are resident in and in the country which is designated for dispute resolution over the \textsc{agreement} and its execution. 
        \subsection{Material copy}
        In the event of technical or legal inability to apply the digital signature, the \textsc{agreement} shall be produced on paper in two copies and signed on each page by the \textsc{customer} on the left and by the \textsc{provider} on the right. The \textsc{agreement} is considered in force after the moment both of the \textsc{parties} sign it.}
        {}
      \clause
        {Строк дії}
        {Duration}
        {Eesti}
        {У\fauxsc{года} укладається \fauxsc{сторонами} із строком дії у \fieldw{duration}{75} рік(-и,-ів). Якщо \fauxsc{замовник} на дату закінчення дії \fauxsc{угоди} продовжує потребувати послуг від \fauxsc{виконавця}, \fauxsc{сторони} мають укласти нову \fauxsc{угоду}.}
        {The \textsc{agreement} is concluded by the \textsc{parties} for the duration of \fieldw{duration}{75} year(s). Should the \textsc{customer} require further service on the \textsc{provider}'s behalf, a new \textsc{agreement} \emph{shall} be concluded.}
        {}
      \clause
        {Припинення дії}
        {Termination}
        {Eeesti}
        {Дія \fauxsc{угоди} припиняється у наступних випадках:
        \begin{enumerate}
          \item Закінчення строку дії;
          \item Взаємна згода \fauxsc{сторін};
          \item В односторонньому порядку, у випадку чого \fauxsc{угода} діятиме 2 календарні тижні із дати інформування іншої \fauxsc{сторони} про намір розірвати угоду; 
          \item Визнання угоди нечинною через невідповідність законодавству;
          \item Рішення компетентного суду;
        \end{enumerate}}
        {The \textsc{agreement} is terminated under the following circumstances:
        \begin{enumerate}
          \item The duration being expired;
          \item Mutual agreement by the \textsc{parties};
          \item Unilaterally, in this case the \textsc{agreement} will be in force for two weeks from the date when the other \textsc{party} has been informed about the intention to terminate the \textsc{agreement};
          \item Invalidation of the \textsc{agreement} in case of legislatory non-compliance;
          \item The decision of the competent court; 
        \end{enumerate} }
        {}
      \clausenewpage
        {Підписанти}
        {Signatories}
        {}
        {\textbf{ЗАМОВНИК}\\
        \fieldline{customernameUK}\\
        Код (податковий номер):\\
        \fieldline{customercodeUKETET}\\
        Країна:\\
        \fieldline{customercountryUK}\\
        Юридична адреса:\\
        \fieldline{customeraddr1UK}\\
        \fieldline{customeraddr2UK}\\
        \\
        Що його представляє:\\
        \fieldline{customerpersonUK}\\
        На підставі\\
        \fieldline{customerrightUK}\\
        Індивідуальний код представника:\\
        \fieldline{customeridUK}\\
        Країна, що надала код:\\
        \fieldline{customerperscountryUK}\\
        \\
        \textbf{ВИКОНАВЕЦЬ}\\
        \fieldline{providernameUK}\\
        Код (податковий номер):\\
        \fieldline{providercodeUKETET}\\
        Країна:\\
        \fieldline{providercountryUK}\\
        Юридична адреса:\\
        \fieldline{provideraddr1UK}\\
        \fieldline{provideraddr2UK}\\
        \\
        Що його представляє:\\
        \fieldline{providerpersonUK}\\
        На підставі\\
        \fieldline{providerrightUK}\\
        Індивідуальний код представника:\\
        \fieldline{provideridUKENET}\\
        Країна, що надала код:\\
        \fieldline{providerperscountryUK}\\
        }
        {\textbf{CUSTOMER}\\
        \fieldline{customernameEN}\\
        Code (tax number):\\
        \fieldline{customercodeUKETET}\\
        Country:\\
        \fieldline{customercountryEN}\\
        Legal address:\\
        \fieldline{customeraddr1EN}\\
        \fieldline{customeraddr2EN}\\
        \\
        Represented by:\\
        \fieldline{customerpersonENET}\\
        On the basis of\\
        \fieldline{customerrightEN}\\
        Representant's individual number:\\
        \fieldline{customeridUKENET}\\
        Country which issued the number:\\
        \fieldline{customerperscountryEN}\\
        \\
        \textbf{PROVIDER}\\
        \fieldline{providernameEN}\\
        Code (tax number):\\
        \fieldline{providercodeUKETET}\\
        Country:\\
        \fieldline{providercountryEN}\\
        Legal address:\\
        \fieldline{provideraddr1EN}\\
        \fieldline{provideraddr2EN}\\
        \\
        Represented by:\\
        \fieldline{providerpersonENET}\\
        On the basis of\\
        \fieldline{providerrightEN}\\
        Representant's individual number:\\
        \fieldline{provideridUKENET}\\
        Country which issued the number:\\
        \fieldline{providerperscountryEN}\\}
        {\textbf{@}\\
        \fieldline{customernameEN}\\
        @:\\
        \fieldline{customercodeUKETET}\\
        @:\\
        \fieldline{customercountryEN}\\
        @:\\
        \fieldline{customeraddr1EN}\\
        \fieldline{customeraddr2EN}\\
        \\
        @:\\
        \fieldline{customerpersonENET}\\
        @\\
        \fieldline{customerrightET}\\
        @\\
        \fieldline{customeridUKENET}\\
        @:\\
        \fieldline{customerperscountryET}\\
        \\
        \textbf{@}\\
        \fieldline{providernameET}\\
        @:\\
        \fieldline{providercodeUKETET}\\
        @:\\
        \fieldline{providercountryET}\\
        @:\\
        \fieldline{provideraddr1ET}\\
        \fieldline{provideraddr2ET}\\
        \\
        @:\\
        \fieldline{providerpersonENET}\\
        @\\
        \fieldline{providerrightET}\\
        @:\\
        \fieldline{provideridUKENET}\\
        @:\\
        \fieldline{providerperscountryET}\\}
    \end{paracol}
  \appendix
  %
% Copyright 2016 Erint Labs OÜ
%
% Licensed under the EUPL, Version 1.1 or – as soon they
% will be approved by the European Commission - subsequent
% versions of the EUPL (the "Licence");
% You may not use this work except in compliance with the
% Licence.
% You may obtain a copy of the Licence at:
%
% https://joinup.ec.europa.eu/software/page/eupl
%
% Unless required by applicable law or agreed to in
% writing, software distributed under the Licence is
% distributed on an "AS IS" basis,
% WITHOUT WARRANTIES OR CONDITIONS OF ANY KIND, either
% express or implied.
% See the Licence for the specific language governing
% permissions and limitations under the Licence.
%
\pagebreak
\section{Таблиця відповідності організаційно-правових форм господарювання\\Business entity type correspondence table\\Majandusüksuse tüübi vastavus esitatud}
\label{app:correspondence}
\begin{tabular}{ | l | l | l |}
  \hline
    \textbf{Українська} & \textbf{English} & \textbf{Eesti} \\
    \hline
    \textbf{ФОП}, Фізична Особа-Підприємець & \textbf{SP}, Sole Proprietorship & \textbf{FIE}, Füüsilisest Isikust Ettevõtja \\
    \textbf{ТОВ}, Товариство із Обмеженною Відповідальністю & \textbf{LLC}, Private Limited Company & \textbf{OÜ}, Osaühing \\
  \hline
\end{tabular}

  %
% Copyright 2016 Erint Labs OÜ
%
% Licensed under the EUPL, Version 1.1 or – as soon they
% will be approved by the European Commission - subsequent
% versions of the EUPL (the "Licence");
% You may not use this work except in compliance with the
% Licence.
% You may obtain a copy of the Licence at:
%
% https://joinup.ec.europa.eu/software/page/eupl
%
% Unless required by applicable law or agreed to in
% writing, software distributed under the Licence is
% distributed on an "AS IS" basis,
% WITHOUT WARRANTIES OR CONDITIONS OF ANY KIND, either
% express or implied.
% See the Licence for the specific language governing
% permissions and limitations under the Licence.
%
\pagebreak
% This is lame, but I have no will to make it better
\newcommand{\myloop}[3]{\stepcounter{#1}#3\ifthenelse{\value{#1} < \value{#2}}{\\\myloop{#1}{#2}{#3}}{\\}}
\section{Список послуг, що надаватимуться\\List of services to be provided\\Loetelu osutatavate teenuste}
\label{app:services}
\newcounter{n}
\setcounter{n}{20}
\newcounter{i}
\begin{tabular}{ | r | l | l | l | l | }
  \hline
    & & & \textbf{Опис послуги} & \textbf{Ціна, €/г} \\
    & \textbf{ДК 009:2010} & \textbf{EMTAK 2008} & \textbf{Service description} & \textbf{Price, €/h} \\
    № & \textbf{КВЕД 2010} & \textbf{NACE Rev.2} & \textbf{Ekspluatatsioonikirjeldust} & \textbf{Hind, €/t} \\
  \hline
    \setcounter{i}{0}
    \myloop{i}{n}{\arabic{i} &
      \raisebox{-2.2pt}{\TextField[borderwidth=0,width=40pt,height=10pt,charsize=8pt,name=kved\arabic{i}]{}} &
      \raisebox{-2.2pt}{\TextField[borderwidth=0,width=40pt,height=10pt,charsize=8pt,name=emtak\arabic{i}]{}} &
      \raisebox{-2.2pt}{\TextField[borderwidth=0,width=560pt,height=10pt,charsize=8pt,name=desc\arabic{i}]{}} &
      \raisebox{-2.2pt}{\TextField[borderwidth=0,width=40pt,height=10pt,charsize=8pt,name=rate\arabic{i}]{}}}
    \hline
\end{tabular}

  %
% Copyright 2016 Erint Labs OÜ
%
% Licensed under the EUPL, Version 1.1 or – as soon they
% will be approved by the European Commission - subsequent
% versions of the EUPL (the "Licence");
% You may not use this work except in compliance with the
% Licence.
% You may obtain a copy of the Licence at:
%
% https://joinup.ec.europa.eu/software/page/eupl
%
% Unless required by applicable law or agreed to in
% writing, software distributed under the Licence is
% distributed on an "AS IS" basis,
% WITHOUT WARRANTIES OR CONDITIONS OF ANY KIND, either
% express or implied.
% See the Licence for the specific language governing
% permissions and limitations under the Licence.
%
\pagebreak
\section{Банківські реквизити сторін\\Bank details of the parties\\Pangaandmed poolte}
\label{app:bank}
\newcommand\banktable[1]{
\begin{tabular}{ | r | l | }
  \hline
  Отримувач & \multirow{3}{*}{\fieldb{benef#1}} \\
  Beneficiary & \\
  Kasusaaja & \\
  \hline
  Банк отримувача & \multirow{3}{*}{\fieldb{bank#1}} \\
  Beneficiary bank & \\
  Saaja pank & \\
  \hline
  Відділення банку отримувача & \multirow{3}{*}{\fieldb{branch#1}} \\
  Beneficiary bank branch & \\
  Saaja panga filiaal & \\
  \hline
  IBAN & \multirow{4}{*}{\fieldbs{iban#1}} \\
  Номе рахунку отримувача & \\
  Beneficiary account \# & \\
  Saaja konto \# & \\
  \hline
  BIC/SWIFT банку отримувача & \multirow{3}{*}{\fieldb{swift#1}} \\
  Beneficiary bank BIC/SWIFT & \\
  Saaja panga BIC/SWIFT & \\
  \hline
  Банк-кореспонтент & \multirow{3}{*}{\fieldb{corrbank#1}} \\
  Correspondent bank & \\
  Korrespondentpank & \\
  \hline
  Номер кореспондентського рахунку & \multirow{3}{*}{\fieldbs{corriban#1}} \\
  Correspondent account \# & \\
  Korrespondentkonto \# & \\
  \hline
  BIC/SWIFT банку-кореспонденту & \multirow{3}{*}{\fieldb{corrswift#1}} \\
  Correspondent bank BIC/SWIFT & \\
  Korrespondentpangale BIC/SWIFT & \\
  \hline
  Розташування банку-кореспонденту & \multirow{3}{*}{\fieldb{corr#1}} \\
  Correspondent bank location & \\
  Korrespondentpanga asukohta & \\
  \hline
\end{tabular}
}
\begin{paracol}{2}
  \section*{ЗАМОВНИК / PROVIDER / KLIENT}
  \switchcolumn
  \section*{ВИКОНАВЕЦЬ / PROVIDER / TÖÖVÕTJA}
  \switchcolumn*
  \banktable{customer}
  \switchcolumn
  \banktable{provider}
\end{paracol}

  \end{Form}
\end{document}
